\documentclass[a4paper,12pt]{article}
\usepackage[utf8]{inputenc}
\usepackage[italian]{babel}
\usepackage{listings}
\usepackage{xcolor}
\usepackage{float}
\usepackage{placeins}
\usepackage{url}
\setlength{\emergencystretch}{3em}

\lstset{
  language=C,
  backgroundcolor=\color{lightgray!20},
  frame=single,
  breaklines=true,
  numbers=left,
  numberstyle=\tiny,
  stepnumber=1,
  numbersep=5pt,
  captionpos=b,
  basicstyle=\ttfamily\small,
  keywordstyle=\color{blue}\bfseries,
  commentstyle=\color{green!60!black},
  stringstyle=\color{red},
  showstringspaces=false
}

\title{LabReSiD26 -- Hands-On 14}
\author{Davide Ferrara -- 518629}
\date{}

\begin{document}

\maketitle

% ---------------------------------------------------------------------------
\section{Esercizio 1: Conto corrente condiviso}
% ---------------------------------------------------------------------------

$N = 4$ sportelli operano in parallelo sullo stesso conto corrente condiviso,
eseguendo in loop operazioni di deposito e prelievo.
Il prelievo è condizionato al saldo disponibile, che non deve mai scendere
sotto zero.

\subsection*{Race condition senza mutex}

Il controllo \texttt{if (saldo >= importo)} da solo non è sufficiente.
Se due thread leggono \texttt{saldo} contemporaneamente, entrambi trovano
la condizione vera ed eseguono entrambi il prelievo: il saldo può diventare
negativo e il risultato cambia ad ogni esecuzione.
L'intera sequenza \emph{leggi / controlla / modifica} deve essere atomica.

\subsection*{Sezione critica con mutex}

Si dichiara un \texttt{pthread\_mutex\_t} globale inizializzato staticamente
con \texttt{PTHREAD\_MUTEX\_INITIALIZER}.
Il mutex protegge sia il controllo \texttt{if} che la modifica di
\texttt{saldo}: senza proteggere entrambi la race condition persiste.

\begin{lstlisting}[caption={Sezione critica protetta da mutex}]
pthread_mutex_t mutex = PTHREAD_MUTEX_INITIALIZER;

/*
 * Per N_OPERAZIONI si calcola un importo random che va a 0 a 100.
 * Se `i` e' pari allora l'importo si somma altrimenti si sottrae.
 * Il controllo if (saldo >= importo) non e' sufficiente in quanto la
 * funzione sportello viene eseguita da piu thread e saldo e' una
 * variabile globale condivisa: se due thread scrivono parallelamente
 * uno potrebbe sovrascrivere il risultato dell'altro.
 * Il MUTEX garantisce che solo un thread alla volta acceda alla
 * sezione critica.
 * */
void *sportello(void *arg) {
  (void)arg;
  unsigned int seed = 518629;
  for (int i = 0; i < N_OPERAZIONI; i++) {
    int importo = (rand_r(&seed) % 100) + 1;

    // Sezione critica
    pthread_mutex_lock(&mutex);
    if (i % 2 == 0) {
      saldo += importo;         /* deposito  */
    } else {
      if (saldo >= importo)     /* prelievo  */
        saldo -= importo;
    }
    pthread_mutex_unlock(&mutex);
    // Fine sezione critica
  }
  return NULL;
}
\end{lstlisting}

\subsection*{Riproducibilità con \texttt{rand\_r}}

\texttt{rand()} usa uno stato interno globale non thread-safe: più thread
che lo chiamano contemporaneamente generano una race condition sullo stato
interno, producendo sequenze diverse ad ogni esecuzione.
Si usa \texttt{rand\_r(\&seed)} con un seed locale per thread, dichiarato
fuori dal loop così che lo stato si aggiorni tra un'iterazione e l'altra.

\begin{lstlisting}[language=bash, caption={Test Esercizio 1 -- risultato riproducibile}]
$ ./banca
Saldo finale: 131056 euro
$ ./banca
Saldo finale: 131056 euro
$ ./banca
Saldo finale: 131056 euro
\end{lstlisting}

% ---------------------------------------------------------------------------
\section{Esercizio 2: Stampante condivisa con timeout}
% ---------------------------------------------------------------------------

Un ufficio ha una sola stampante condivisa tra $N = 6$ utenti.
Ogni utente tenta di acquisire la stampante: se non si libera entro
3 secondi rinuncia. Una stampa richiede 2 secondi e gli utenti
arrivano uno ogni 0.2 secondi.

\subsection*{Timeout con \texttt{pthread\_mutex\_timedlock}}

\texttt{pthread\_mutex\_timedlock()} richiede un tempo \textbf{assoluto}
basato su \texttt{CLOCK\_REALTIME}. Si legge prima l'ora corrente con
\texttt{clock\_gettime()} e poi si aggiungono i secondi di attesa.
Le funzioni \texttt{pthread\_*} restituiscono il codice di errore
direttamente come valore di ritorno, senza settare \texttt{errno}.
Se restituisce \texttt{ETIMEDOUT} il mutex \textbf{non} è stato acquisito e
chiamare \texttt{unlock()} in questo caso produrrebbe comportamento indefinito.

\begin{lstlisting}[caption={Funzione utente con timeout}]
void *utente(void *arg) {
  int id = *(int *)arg;

  /* Ora di sistema attuale (CLOCK_REALTIME) + 3 secondi */
  struct timespec ts;
  clock_gettime(CLOCK_REALTIME, &ts);
  ts.tv_sec += 3;

  /* Provo ad acquisire il mutex per 3 secondi */
  int rc = pthread_mutex_timedlock(&stampante, &ts);
  if (rc == 0) {
    printf("[Utente %d] Sta stampando... \n", id);
    sleep(2); /* Simula la stampa */
    printf("[Utente %d] Ha terminato la stampa :D \n", id);
    pthread_mutex_unlock(&stampante);
  } else if (rc == ETIMEDOUT) {
    printf("[Utente %d] Rinuncia alla stampa :(\n", id);
  } else {
    printf("[Utente %d] Errore nell'acquisone del mutex.\n", id);
  }
  return NULL;
}
\end{lstlisting}

\subsection*{Analisi del comportamento}

Gli utenti arrivano ogni 0.2\,s, ogni stampa dura 2\,s e il timeout è
3\,s. La timeline permette di ragionare su chi riesce a stampare:

\begin{lstlisting}[language=bash]
/*
 *  0, 0.2, 0.4, 0.6, 0.8, 1, 1.2, 1.4, 1.6, 1.8, 2, 2.2, ...
 * u0(stampa) -> sleep(2) .................. u0(termina)
 *     u1(attende) -> wait(3) .............. u1(stampa) -> ...
 *           u2(attende) -> wait(3) ..................... u2(rinuncia)
 * */
\end{lstlisting}

\begin{lstlisting}[language=bash, caption={Test Esercizio 2 -- stampa 2s\, timeout 3s}]
$ ./stampante
[Utente 0] Sta stampando...
[Utente 0] Ha terminato la stampa :D
[Utente 1] Sta stampando...
[Utente 2] Rinuncia alla stampa :(
[Utente 3] Rinuncia alla stampa :(
[Utente 4] Rinuncia alla stampa :(
[Utente 1] Ha terminato la stampa :D
[Utente 5] Sta stampando...
[Utente 5] Ha terminato la stampa :D
\end{lstlisting}

% ---------------------------------------------------------------------------
\section{Esercizio 3: Rubrica telefonica con rwlock}
% ---------------------------------------------------------------------------

Una rubrica è condivisa tra $N = 8$ thread lettori e un thread scrittore.
I lettori fanno ricerche frequenti e possono leggere in parallelo tra loro;
lo scrittore aggiorna i contatti raramente ma richiede accesso esclusivo.
Un mutex ordinario serializzerebbe anche le letture tra loro, inutilmente.
\texttt{pthread\_rwlock\_t} risolve questo: più lettori possono tenere il
lock simultaneamente, mentre lo scrittore ottiene accesso esclusivo.

\subsection*{Lettori in parallelo}

Ogni lettore acquisisce il lock in lettura con
\texttt{pthread\_rwlock\_rdlock()}, legge il contatto e rilascia.
Più lettori possono tenere il lock contemporaneamente, perche leggere
in parallelo è sempre safe.

\begin{lstlisting}[caption={Thread lettore}]
void *lettore(void *arg) {
  int id = *(int *)arg;
  for (int i = 0; i < 5; i++) {
    // Acquisisce il lock in lettura
    pthread_rwlock_rdlock(&rubrica_lock);

    int idx = i % n_contatti;
    printf("Lettore %d: %s -> %s\n", id, rubrica[idx].nome,
           rubrica[idx].numero);
    usleep(100000);

    // Rilascia il lock
    pthread_rwlock_unlock(&rubrica_lock);
  }
  return NULL;
}
\end{lstlisting}

\subsection*{Scrittore con accesso esclusivo}

Lo scrittore dorme 1 secondo per lasciare partire i lettori, poi acquisisce
il lock in scrittura con \texttt{pthread\_rwlock\_wrlock()}: blocca finché
tutti i lettori attivi non rilasciano il lock, poi ottiene accesso esclusivo 
(scrive solamente quando nessuno legge).

\begin{lstlisting}[caption={Thread scrittore}]
void *scrittore(void *arg) {
  (void)arg;
  sleep(1); /* lascia partire i lettori */

  pthread_rwlock_wrlock(&rubrica_lock);

  printf("Scrittore: aggiorno il numero di Bob...\n");
  strcpy(rubrica[1].numero, "333-9999");
  sleep(2); /* simula aggiornamento lento */
  printf("Scrittore: aggiornamento completato.\n");

  pthread_rwlock_unlock(&rubrica_lock);
  return NULL;
}
\end{lstlisting}

\begin{lstlisting}[language=bash, caption={Test Esercizio 3 -- interleaving, blocco scrittore e aggiornamento}]
$ ./rubrica
Lettore 0: Alice -> 333-1111          # prima tornata: lettori in parallelo
Lettore 3: Alice -> 333-1111          # (interleaving visibile negli id)
...
Lettore 0: Bob -> 333-2222
Lettore 4: Bob -> 333-2222
Lettore 2: Bob -> 333-2222
...
Scrittore: aggiorno il numero di Bob...
                                       # blocco silenzioso 2s (wrlock esclusivo)
Scrittore: aggiornamento completato.
Lettore 0: Alice -> 333-1111          # seconda tornata: lettori ripartono
Lettore 1: Alice -> 333-1111
...
Lettore 0: Bob -> 333-9999            # numero aggiornato visibile a tutti
Lettore 3: Bob -> 333-9999
Lettore 1: Bob -> 333-9999
...
\end{lstlisting}

L'interleaving degli id nei round di lettura conferma il parallelismo dei
lettori. Il blocco silenzioso di 2 secondi tra i messaggi dello scrittore
dimostra l'accesso esclusivo del wrlock. Nella seconda tornata tutti i
lettori leggono \texttt{333-9999}: nessuno vede il vecchio valore dopo
l'aggiornamento. Inoltre il programma non segnala nessuna race condition 
che sarebbe comparsa compilando con \texttt{-fsanitize=thread}.

\end{document}
